\section{Tubular neighbourhoods of Pfaffian hypersurfaces}\label{sec:tubular}
For $A\subseteq\R^n$ and $\varepsilon>0$, write
\[
T(A,\varepsilon)=\{z\in\R^n\colon\dist(z,A)\le\varepsilon\}
\]
for the full closed distance neighbourhood. If $M$ is a smooth $m$-dimensional submanifold of $\R^n$, with the induced Euclidean metric and positive codimension $c=n-m$, define its \emph{normal tube} by
\[
T^\perp(M,\varepsilon)
=\{x+u\colon x\in M,\ u\perp T_xM,\ \|u\|\le\varepsilon\}.
\]
Always $T^\perp(M,\varepsilon)\subseteq T(M,\varepsilon)$. If $M$ is closed as a subset of $\R^n$ and has no boundary, equality holds for every $\varepsilon>0$: a nearest point exists and the displacement to it is normal to $M$. Neither boundedness nor uniqueness of the nearest point is needed. For manifolds with boundary, the two neighbourhoods can differ.

We combine a tube-volume estimate from~\cite{lotz2015volume} with effective bounds for the generalised Gauss map. A localisation argument will allow the hypersurface itself to be unbounded.

\subsection{Generalised Gauss map}
For $x\in M$, write $N_xM=(T_xM)^\perp$ and
\[
S(NM)=\{(x,v)\in M\times S^{n-1}\colon v\in N_xM\}
\]
for the normal space and normal sphere bundle, respectively. The latter is a smooth manifold of dimension $n-1$.

\begin{definition}
The \emph{generalised Gauss map} is
\[
\gamma_M\colon S(NM)\longrightarrow S^{n-1},\qquad (x,v)\longmapsto v.
\]
\end{definition}

For a smooth map $\psi$ between manifolds of the same dimension, define its \emph{maximal degree} by
\[
\md(\psi):=\sup_{v\in\reg(\psi)}|\psi^{-1}(v)|\in\N\cup\{\infty\},
\]
where $\reg(\psi)$ is the set of regular values and $|\psi^{-1}(v)|$ counts distinct points, without signs or multiplicities. Each regular fibre is discrete, and is finite when the domain is compact~\cite[\S 1]{milnor1997topology}. A uniform bound on these cardinalities requires further information; for the Pfaffian manifolds below it follows from effective zero counts.

Let $\cE_j^n$ denote the space of $j$-dimensional affine subspaces of $\R^n$. Almost every $H\in\cE_{c+i}^n$ meets $M$ transversally, and then $M\cap H$ is an $i$-dimensional submanifold of $H\cong\R^{c+i}$. Define
\[
\md(M,H):=\md(\gamma_{M\cap H}),
\qquad
\md_i(M):=\sup_{H\in\cE_{c+i}^n}\md(M,H),
\quad 0\le i\le m,
\]
with $\md(M,H)=0$ for a non-transverse intersection or an empty section. The normal spaces of $M\cap H$ are taken within $H$. We abbreviate $\md(M):=\md_m(M)=\md(\gamma_M)$.

\begin{lemma}\label{lem:degree-restriction}
If $U$ is an open subset of $M$, then $\md_i(U)\le\md_i(M)$ for every $0\le i\le m$.
\end{lemma}

\begin{proof}
Fix a transverse section of $U$, a regular value of its Gauss map, and finitely many points of that fibre. Transversality and the inverse function theorem imply that these distinct regular preimages persist under sufficiently small changes of the affine section and normal direction. Choose a nearby section transverse to all of $M$, and then a nearby direction regular for its Gauss map; such choices exist by transversality and Sard's theorem. The chosen finite collection has cardinality at most $\md_i(M)$. Taking suprema over finite collections, regular values and sections proves the claim.
\end{proof}

The following normal-tube estimate follows from~\cite[Theorems 3.1 and 4.3]{lotz2015volume}, with the boundary convention explained below.

\begin{theorem}[Normal-tube bound]\label{thm:tube_bound}
Let $M\subset B(p,\rho)$ be a compact smooth submanifold of dimension $0\le m<n$, possibly with boundary, and set $c=n-m$. Then
\[
\vol T^\perp(M,\varepsilon)\le
2\omega_n\rho^n\sum_{i=0}^m\binom n{c+i}\md_i(M^\circ)
\left(\frac{\varepsilon}{\rho}\right)^{c+i},
\]
where $M^\circ=M\setminus\partial M$ and $\varepsilon,\rho>0$.
\end{theorem}

\begin{remark}\label{rem:tube_bound_boundary}
Theorem 3.1 of~\cite{lotz2015volume} bounds the normal tube based on an open subset $U\subseteq M\setminus\partial M$. Apply it with $U=M^\circ$ and use the section-degree argument of its Theorem 4.3; the extension of the Gauss-map estimate to such $U$ is stated after its Lemma 4.2. This gives the displayed bound for $\vol T^\perp(M^\circ,\varepsilon)$. Normal vectors of length at most $\varepsilon$ based on $\partial M$, with normal taken relative to $M$, form a disk bundle of dimension $(m-1)+c=n-1$. Its image has zero $n$-dimensional volume, so $\vol T^\perp(M,\varepsilon)=\vol T^\perp(M^\circ,\varepsilon)$.

This does not bound the full distance neighbourhood $T(M,\varepsilon)$ when $M$ has boundary. For example, the normal tube of a line segment in $\R^2$ is a rectangle, whereas its distance neighbourhood also includes two semicircular caps. The tube $T^\perp(\partial M,\varepsilon)$ uses normals to $\partial M$ as a separate submanifold and can have positive $n$-dimensional volume.
\end{remark}

\subsection{Localisation and Gauss-map degree bounds}
The next observation avoids introducing a separate tube around the boundary of a truncation.

\begin{proposition}\label{prop:local-tube}
Let $V\subset\R^n$ be a closed smooth hypersurface without boundary, with finite degrees $\md_i(V)$. If $X$ is uniform in $B(p,\rho)$, where $\rho,\varepsilon>0$, and $u=\varepsilon/\rho$, then
\begin{equation}\label{eq:local-tube}
\prob\{\dist(X,V)\le\varepsilon\}\le
2\sum_{i=0}^{n-1}\binom n{i+1}\md_i(V) u^{i+1}(1+u)^{n-i-1}.
\end{equation}
\end{proposition}

\begin{proof}
The empty case is immediate. Choose $R>\rho+\varepsilon$ such that the sphere $S^{n-1}(p,R)$ meets $V$ transversally. Then $M_R=V\cap B(p,R)$ is a compact manifold with boundary. If $z\in B(p,\rho)$ has distance at most $\varepsilon$ from $V$, choose a nearest point $x\in V$. We have
\[
\|x-p\|\le\|x-z\|+\|z-p\|\le\varepsilon+\rho<R,
\]
so $x\in M_R^\circ$. The nearest-point condition gives $z-x\perp T_xV$, and hence
\[
T(V,\varepsilon)\cap B(p,\rho)
\subseteq T^\perp(M_R^\circ,\varepsilon)
\subseteq T^\perp(M_R,\varepsilon).
\]
Apply Theorem~\ref{thm:tube_bound} to the compact manifold $M_R$, and Lemma~\ref{lem:degree-restriction} to its interior, which is open in $V$. Dividing by $\omega_n\rho^n$ bounds the probability by
\[
2\sum_{i=0}^{n-1}\binom n{i+1}\md_i(V)
u^{i+1}\left(\frac R\rho\right)^{n-i-1}.
\]
By Sard's theorem, regular radii can be chosen decreasing to $\rho+\varepsilon$. Taking this limit proves~\eqref{eq:local-tube}.
\end{proof}

For hypersurfaces, a Gauss fibre is described by a square system without auxiliary variables. The following non-degeneracy statement will be used in both the Pfaffian and Bernstein bounds. In what follows, the notation $\partial_{\tau}f=\langle\nabla f,\tau\rangle$ is used for the directional derivative.

\begin{lemma}\label{le:gauss-nondegenerate}
Let $f:\R^d\to\R$ be smooth with $\nabla f\ne0$ on $V=\mathcal Z(f)$, and let $v\in S^{d-1}$ be a regular value of $\gamma_V$. Choose any basis $\tau_1,\dots,\tau_{d-1}$ of $v^\perp$. Then the solutions of
\begin{equation}\label{eq:gauss-directional-system}
f=\partial_{\tau_1}f=\cdots=\partial_{\tau_{d-1}}f=0
\end{equation}
are in bijection with $\gamma_V^{-1}(v)$, and every solution is non-degenerate. In particular, when $v_d\ne0$, choosing $\tau_j=v_de_j-v_je_d$ gives
\begin{equation}\label{eq:gauss-cross-system}
f=0,\qquad v_d\partial_j f-v_j\partial_d f=0,\quad 1\le j<d.
\end{equation}
\end{lemma}

\begin{proof}
The derivative equations say exactly that the nonzero gradient of $f$ is parallel to $v$. Each such point $x$ contributes precisely the bundle point $(x,v)$ to the specified fibre.

A rotation and an invertible change of the derivative equations reduce non-degeneracy to $v=e_d$ and $\tau_j=e_j$. At a solution $x$, we have $\partial_d f(x)\ne0$, and the implicit function theorem represents $V$ locally as $u\mapsto(u,h(u))$. Write $x=(u_0,h(u_0))$; then $\nabla h(u_0)=0$. Twice differentiating $f(u,h(u))=0$ gives
\[
\partial_{jk}f(x)=-\partial_d f(x)\,\partial_{jk}h(u_0)
\quad(1\le j,k<d).
\]
Regularity of the Gauss map is equivalent to nonsingularity of $\operatorname{Hess}h(u_0)$. Expanding the Jacobian determinant of $(f,\partial_1f,\dots,\partial_{d-1}f)$ along its first row gives absolute value
\[
|\partial_d f(x)|^d\,|\det\operatorname{Hess}h(u_0)|>0.
\]
For $d=1$, the system is simply $f=0$ and the assertion follows from $f'(x)\ne0$.
\end{proof}

\begin{proposition}\label{prop:hypersurface-degree}
Let $f\in\Pfaff_{\alpha,\beta,s}(\R^n)$, where $\alpha,\beta\ge1$, and assume that $0$ is a regular value. For $V=\mathcal Z(f)$ and $0\le i<n$,
\begin{equation}\label{eq:hypersurface-ideg}
\md_i(V)\le
2^{\frac{s(s-1)}2}\beta(\alpha+\beta-1)^i
\left[\beta+i(\alpha+\beta-2)+\min\{i+1,s\}\alpha\right]^s.
\end{equation}
\end{proposition}

\begin{proof}
Parametrise a transverse affine section $H\in\cE_{i+1}^n$ by an affine isometry from $\R^{i+1}$. By Remark~\ref{re:affine}, the restricted function has the same format and $0$ remains a regular value. For each regular Gauss direction, Lemma~\ref{le:gauss-nondegenerate} supplies a non-degenerate square system with one equation of degree $\beta$ and $i$ equations of degree at most $\alpha+\beta-1$, by Lemma~\ref{le:pfaffalgebra}(3). Theorem~\ref{thm:khovanskii} gives~\eqref{eq:hypersurface-ideg}. The bound holds for every such fibre, so taking the two suprema proves the result without a compactness assumption.
\end{proof}

For completeness, a similar argument with Lagrange multipliers gives a bound for complete intersections.

\begin{proposition}\label{prop:pfaffian-degree-bound}
Let $M=\mathcal Z(f_1,\dots,f_c)\subset\R^n$, where $1\le c\le n$, the gradients $\nabla f_j$ are linearly independent on $M$, and the functions share a Pfaffian chain with formats $(\alpha,\beta_j,s)$, $\alpha,\beta_j\ge1$. Put $\overline\beta=\max_j\beta_j$. Then
\begin{align*}
\md(M)\le{}&
2^{\frac{s(s-1)}2}(\alpha+\overline\beta)^n\prod_{j=1}^c\beta_j\\
&{}\cdot\left[
n(\alpha+\overline\beta-1)+\sum_{j=1}^c(\beta_j-1)
+\min\{n+c,s\}\alpha+1
\right]^s.
\end{align*}
The same bound in ambient dimension $c+i$ bounds $\md_i(M)$.
\end{proposition}

\begin{proof}
For a regular value $v$ of $\gamma_M$, its fibre is in bijection with solutions of
\begin{equation}\label{eq:pfaffian-system}
f_j(x)=0\quad(1\le j\le c),\qquad
\sum_{j=1}^c\lambda_j\nabla f_j(x)=v.
\end{equation}
The multipliers are unique by independence of the gradients. To check non-degeneracy, identify $(x,\lambda)\in M\times\R^c$ with $(x,\sum_j\lambda_j\nabla f_j(x))$ in the normal bundle. Near a solution, projection of this bundle to $\R^n$ is, in radial coordinates, $(x,u,r)\mapsto r\gamma_M(x,u)$. Its differential is invertible at $r=1$ precisely when $\diff\gamma_M$ is invertible. Since $(f_1,\dots,f_c)$ is a submersion on $M$, this is equivalent to the full Jacobian of~\eqref{eq:pfaffian-system} being nonsingular. Theorem~\ref{thm:khovanskii} now applies in $n+c$ variables, with degrees $\beta_1,\dots,\beta_c$ and $n$ degrees at most $\alpha+\overline\beta$. Restriction to a transverse affine section preserves the formats.
\end{proof}

\subsection{A Pfaffian tube formula}
Combining the local tube bound with the direct hypersurface system gives the following estimate.

\begin{theorem}\label{thm:prob_bound}
Let $f\in\Pfaff_{\alpha,\beta,s}(\R^n)$ with $\alpha,\beta\ge1$ and $0$ a regular value, and set $V=\mathcal Z(f)$. Let $X$ be uniformly distributed in $B(p,\rho)$, where $p\in\R^n$ and $\rho>0$. For every $\varepsilon>0$,
\[
\prob\{\dist(X,V)\le\varepsilon\}\le
C_{\alpha,\beta,s,n}
\left[\left(1+(\alpha+\beta)\frac{\varepsilon}{\rho}\right)^n
-\left(1+\frac{\varepsilon}{\rho}\right)^n\right],
\]
where
\begin{equation}\label{eq:pfaffian-tube-constant}
C_{\alpha,\beta,s,n}
=\frac{2^{\frac{s(s-1)}2+1}\beta}{\alpha+\beta-1}
\left[\beta+(n-1)(\alpha+\beta-2)+\min\{n,s\}\alpha\right]^s.
\end{equation}
\end{theorem}

\begin{proof}
Put $D=\alpha+\beta-1$, $B=\beta+(n-1)(\alpha+\beta-2)+\min\{n,s\}\alpha$, and $A=2^{s(s-1)/2}\beta B^s$. Proposition~\ref{prop:hypersurface-degree} gives $\md_i(V)\le AD^i$. The zero set $V$ is closed, so Proposition~\ref{prop:local-tube} applies. With $u=\varepsilon/\rho$, the binomial identity
\[
\sum_{i=0}^{n-1}\binom n{i+1}D^i u^{i+1}(1+u)^{n-i-1}
=\frac{(1+(D+1)u)^n-(1+u)^n}{D}
\]
gives the result, since $C_{\alpha,\beta,s,n}=2A/D$.
\end{proof}

\begin{remark}\label{rem:prob_bound_simple}
The right-hand side vanishes with $\varepsilon$ and has expansion
\[
C_{\alpha,\beta,s,n}\,n(\alpha+\beta-1)\frac{\varepsilon}{\rho}
+O_{\alpha,\beta,s,n}\!\left(\frac{\varepsilon^2}{\rho^2}\right)
\quad\text{as }\varepsilon/\rho\to0.
\]
This linear leading term yields the $O(1/t)$ condition-number tails in Section~\ref{sec:robustness}. Every probability bound may also be replaced by its minimum with $1$.
\end{remark}
